\subsection{某些 Shimura 簇模 $p$ 纤维上的 Tate 轮}
我们运用上述机制，对许多 Shimura 簇的模 $p$ 纤维验证 Tate 猜想的「一般」情形。

设 $(G,X,K)$ 如上一小节，$p>2$ 为非分歧素数。
设 $(\hat G, \hat B,\hat T)$ 是 $G_{\bQ_p}$ 的 Langlands 对偶群，配有注 \ref{R: birth of the dual group} 中的 Borel 子群与极大环面。如前固定同构 $\iota: \overline\bQ_p\simeq \bC$，并把 Shimura 余特征标 $\mu$ 看作 $(\hat G,\hat B,\hat T)$ 的支配权。设 $v\mid p$ 是由 $\iota$ 确定的 $E$ 的素位，$\overline{\scrS}$ 如前为 $\Sh_K(G,X)$ 在 $\mO_{E,(v)}$ 上的典范整模型的模 $p$ 纤维。

对 $\la\in\xch(\hat T)$ 与 $\hat G$ 的表示 $V$，记 $V(\la)$ 为 $V$ 的 $\la$-权子空间（关于 $\hat T$）。
定义如下格
\[\Lambda^\Tate_p=\Big\{\la\in \xch(\hat T)\; \Big|\; \sum_{i=0}^{m-1} \sigma_p^i(\la)\in \xcoch(Z_G)\Big\}\subset \xch(\hat T).\]
这里 $Z_G$ 表示 $G$ 的中心，$\xcoch(Z_G)$ 表示 $Z_G$ 的余权格，它是 $\xch(\hat T)$ 的自然子群。对 $\hat G$ 的表示 $V$，定义空间
\[V^\Tate=\bigoplus_{\la\in\Lambda^\Tate_p} V(\la).\]

我们特别关心条件 $V^\Tate_\mu\neq 0$，其中 $\mu$ 是 $G$ 的 Shimura 余特征标。
如 \cite{XZ} 的引言所述，\zy{[来源：Tate cycle 译作「Tate 轮」，沿用《仿射 Grassmann 簇与几何 Satake 等价》（SY02-17）的术语表；轮类映射即 cycle class map。]}在注 \ref{R: coh of Shimura} 关于 Shimura 簇上同调的猜想性描述之下，若 Hecke 模 $\pi_f$ 在 $p$ 处的 Satake 参数足够一般，则该向量空间维数的某个倍数 $a(\pi_f)$ 应当等于 $\overline{\scrS}$ 中间维数紧支上同调的 $\pi_f$-分支中 Tate 轮空间维数。
此外，该空间通常很大。例如当 $G$ 是二次数域上的奇酉（similitude）群、签名 $(i,n-i)$ 时，该空间在惰性素数处的维数为 $\begin{pmatrix}\frac{n+1}{2}\\ i\end{pmatrix}$。

我们还需要一个记号。对代数闭域上维数为 $d$ 的（不必不可约的）代数簇 $Z$，记 ${\rm H}^{\rm BM}_{2d}(Z)(-d)$ 为最高次 Borel–Moore 同调的 $(-d)$-Tate 扭转，它是由 $Z$ 的不可约分支张成的向量空间。现设 $X$ 是有限域 $\FF_q$ 上维数为 $d+r$ 的光滑簇，$Z \subseteq X_{\overline \FF_q}$ 是维数为 $d$ 的（不必不可约的）射影子簇。则有轮类映射
\[\on{cl}: {\rm H}^{\rm BM}_{2d}(Z)(-d)\to \bigcup_{j\geq 1}{\rm H}_{\et,c}^{2r}(X_{\overline \FF_q}, \bQ_\ell(r))^{\sigma_q^j}.\]

\cite{XZ} 的主要定理如下。$\overline{\mathscr S}$ 上有自然的 Newton 分层。
记 $\overline{\mathscr S}_b$ 为基本 Newton 层，$\mH_K$ 通过（上同调）对应作用在其上（这些对应由上一小节构造的 $S$-算子给出）。
\begin{thm}
\label{T:main theorem}
假设 $(G,X)$ 是 Hodge 型，且 $G$ 的中心 $Z_G$ 连通。设 $K\subset G(\bA_f)$ 是开紧子群，记 $2d=\dim\Sh_K(G,X)$。
设 $p>2$ 是关于 $(G,X,K)$ 的非分歧素数，且 $r:=V_{\mu^*}^\Tate\neq 0$。
则：
\begin{enumerate}
\item $\overline{\scrS}$ 的基本 Newton 层 $\overline{\mathscr S}_b$ 是纯 $d$ 维的。特别地，$d$ 总是整数。此外，存在 $\mH_K$-等变同构
$${\rm H}^{\rm BM}_{2d}(\overline{\mathscr S}_{b, \overline \FF_v})(-d)\cong C(G'(\bQ)\backslash G'(\bA_f)/K)^{\oplus r}.$$
这里 $G'$ 是 $G$ 的内形式，它在 $\bA_f$ 上与 $G$ 同构，且在无穷处模中心紧；$C(G'(\bQ)\backslash G'(\bA_f)/K)$ 表示有限集 $G'(\bQ)\backslash G'(\bA_f)/K$ 上的函数空间。

\item
设 $\pi_f$ 是 $\calH_K$ 的不可约模，记
$${\rm H}^{\rm BM}_{2d}(\overline{\mathscr S}_{b, \overline \FF_v})[\pi_f]=\Hom_{\mH_K}(\pi_f,{\rm H}^{\rm BM}_{2d}(\overline{\mathscr S}_{b, \overline \FF_v})(-d))\otimes \pi_f$$
为 $\pi_f$-等型分支。则轮类映射
\[\on{cl}:{\rm H}^{\rm BM}_{2d}(\overline{\mathscr S}_{b, \overline \FF_v})(-d)\to {\rm H}^{2d}_{\et,c}\big(\overline{\mathscr S}_{\overline \FF_v},\overline\bQ_\ell(d)\big)\]
限制在 ${\rm H}^{\rm BM}_{2d}(\overline{\mathscr S}_{b, \overline \FF_v})[\pi_f]$ 上时，若 $\pi_{f,p}$（即 $\pi_f$ 在 $p$ 处的分量）的 Satake 参数是 $V_\mu$-一般的，则它是单射。

\item 假设 $\Sh_K(G,X)$ 是四元数 Shimura 簇或 Kottwitz 算术簇。则当 $\pi_{f,p}$ 的 Satake 参数强 $V_\mu$-一般时，$\on{cl}$ 的 $\pi_f$-等型分支满射到 $T^d(\pi,\Ql)\otimes\pi_f$。特别地，对这些 $\pi_f$，Tate 猜想成立。
\end{enumerate}
\end{thm}
\begin{rmk}
(1) 中心 $Z_G$ 连通这一假设并非本质，只是使表述更简单。

(2) 对 $\hat G$ 的表示 $V$，「$V$-一般」与「强 $V$-一般」Satake 参数的定义见 \cite[定义 1.4.2]{XZ}。$\hat G\sigma_p$ 中的正则半单元素总是 $V$-一般的，反之不然。见 \cite[注 1.4.3]{XZ}。

(3) 该定理的一些特殊情形最初在 \cite{HTX,TX} 中得到证明。
\end{rmk}

该定理的证明依赖若干不同要素。借助 Shimura 簇基本链的 Rapoport–Zink 均化，第 (1) 部分可归约为某些仿射 Deligne–Lusztig 簇的不可约分支问题，后者已在 \cite[\S 4]{XZ} 中研究。
最困难的部分是 (2)，我们通过计算 $\overline\scrS_b$ 中所有 $d$ 维轮之间的相交数来证明它。由第 (1) 部分，$\overline{\scrS}_b$ 中 $d$ 维轮的相交数可编码为一个 $r\times r$-矩阵，其元素在 $p$ 处的球面 Hecke 代数中。一般来说，直接显式计算这个矩阵似乎没有希望。
不过，与例 \ref{intersection by geom Sat} 中的玩具模型类似，该矩阵可以理解为 $\Coh^{\hat G}_{fr}(\hat G\sigma_p)$ 中某些态射的复合。具体地说，先把 $G'(\bQ)\backslash G'(\bA)/K$ 实现为 Shimura 集，其 Shimura 余特征标记为 $\tau$\footnote{$\tau$ 的选取有一点微妙之处，此处从略，见 \cite[注 7.4.4]{XZ}。}。
然后由定理 \ref{introT: spectral action}，该矩阵可计算为
\[ \Hom_{\Coh^{\hat G}_{fr}(\hat G\sigma_p)}(\widetilde{V_{\tau}},\widetilde{V_{\mu}})\otimes  \Hom_{\Coh^{\hat G}_{fr}(\hat G\sigma_p)}(\widetilde{V_{\mu}},\widetilde{V_{\tau}})\to  \Hom_{\Coh^{\hat G}_{fr}(\hat G\sigma_p)}(\widetilde{V_{\tau}},\widetilde{V_{\tau}})=\mathbf{J}.\]
为证明第 (2) 部分，需要判定该配对何时非退化。这本身是表示论中有趣的问题，其解决依赖于对向量值函数的 Chevalley 限制映射的研究。该矩阵的行列式已在 \cite{XZ2} 中算出，由此得到 (2)。第 (3) 部分通过比较两个迹公式证明：$G$ 的 Lefschetz 迹公式与 $G'$ 的 Arthur–Selberg 迹公式。细节参见 \cite[\S 2]{XZ}。

\begin{ex}
设 $G=\on{GU}(1,2r)$ 是虚二次扩张 $E/\bQ$ 对应的 $(2r+1)$ 个变量的酉相似群，其无穷处签名为 $(1,2r)$。它带有标准的 Shimura 数据。固定层级 $K\subset G(\bA_f)$，记 $\Sh_K(G,X)$ 为相应的 Shimura 簇。特别地当 $r=1$ 时，它就是 Picard 模曲面。取 $p$ 为非分歧惰性素数。此时 $\overline{\scrS}_b$ 是一些 Deligne–Lusztig 簇之并，由 $G'(\bQ)\backslash G'(\bA_f)/K$ 参数化，其中 $G'=\on{GU}(0,2r+1)$ 在所有有限素位与 $G$ 同构。这些轮在 $\overline{\scrS}_b$ 中的相交模式已在 \cite{VW} 中给出。然而这些轮之间的相交数要难算得多。事实上，除把定理 \ref{introT: spectral action} 用在此情形之外，我们并不知道如何对一般的 $r$ 直接计算它们（$r=1$ 的情形可以直接处理）。

此时 $\hat G=\GL_{2r+1}\times \bG_m$，$\sigma_p$ 在其上作用为 $(A,c)\mapsto (J(A^T)^{-1}J,c)$，其中 $J$ 是反对角元素全为 $1$ 的反对角矩阵。表示 $V_\mu$ 是 $\GL_{2r+1}$ 的标准表示，$\bG_m$ 通过数乘作用。可以验证 $\dim V_\mu^{\Tate}=1$（这与上面提到的 $\overline{\scrS}_b$ 的不可约分支由 $G'(\bQ)\backslash G'(\bA_f)/K$ 参数化相符）。
可取 $\tau=((1,\ldots,1),1)$，并把 $G'(\bQ)\backslash G'(\bA_f)/K$ 看作 Shimura 余特征标为 $\tau$ 的 Shimura 集。
这里按惯例把 $\hat G$ 的权格等同于 $\bZ^{2r+1}\times\bZ$。则 $\Hom_{\Coh^{\hat G}_{fr}(\hat G\sigma_p)}(\widetilde{V_{\tau}},\widetilde{V_{\mu}})$ 是 $\mathbf J$ 上的自由秩一 module。取定生成元 $\mathbf a_{\on{in}}$，它诱导 $\mathbf J\times\mH^p$-等变同态
$$S(\mathbf a_{\on{in}}): C(G'(\bQ)\backslash G'(\bA_f)/K)\to \on{H}_c^{2r}(\overline{\scrS}_{\overline{\bF}_p},\Ql(r)),$$
实现 $\overline\scrS_b$ 的轮类映射（相差一个倍数）。模 $\Hom_{\Coh^{\hat G}_{fr}(\hat G\sigma_p)}(\widetilde{V_{\mu}},\widetilde{V_{\tau}})$ 也是 $\mathbf J$ 上的自由秩一模。取定生成元 $\mathbf a_\out$，则复合
\[S(\bba_\out)\circ S(\bba_{\on{in}})=S(\bba_\out\circ\bba_{\on{in}})\]
计算这些轮（由 $\overline\scrS_b$ 的不可约分支给出）的相交矩阵。

元素 $h:=\bba_\out\circ\bba_{\on{in}} \in \mathbf J$ 已在 \cite[例 6.4.2]{XZ2} 中显式计算（差一个显然的修正）。在 Satake 同构 $\mathbf J\cong H_{G_{\bQ_p}}\otimes \Ql$ 下，它也可以显式写成 Hecke 算子的组合（\cite{XZ3}）
\begin{equation}
\label{E:Hecke operator in Tpj}
h=p^{r(r+1)}\sum_{i=0}^{r} (-1)^i(2i+1)p^{(i-r)(r+i+1)}\sum_{j=0}^{r-i} \begin{bmatrix} 2r+1-2j \\ r-i- j\end{bmatrix}_{v=-p}T_{p,j}.
\end{equation}
这里 $T_{p,j}=1_{K_p\la_j(p)K_p}$，其中 $\la_i=(1^i,0^{2r-2i+1},(-1)^i,0)$；$\begin{bmatrix} n \\ m\end{bmatrix}_v$ 是二项式系数的 $v$-类比
\[
[0]_v=1,\quad [n]_v=\frac{v^n-1}{v-1},\quad [n]_v!=[n]_v[n-1]_v\cdots[1]_v,\quad \begin{bmatrix} n \\ m\end{bmatrix}_v=\frac{[n]_v!}{[n-m]_v![m]_v!}.
\]
换言之，此情形下 $\overline\scrS_b$ 中轮的相交矩阵由 Hecke 算子 \eqref{E:Hecke operator in Tpj} 计算。细节与更多讨论参见 \cite{XZ3}。
\end{ex}
